\section*{Finite stock and recovery}
Atoms and molecules are suspended, enlarged teaching models; this deliberate departure from real laboratory behaviour is consistent across tasks.

The shelf begins with \textbf{8 C, 8 O and 32 H}, all individually visible in four shallow baskets. Ordered sockets avoid a physics pile, obscure overlaps and uncontrolled bouncing. Empty sockets and counts make the system state visible. Basket arrangement is staggered front-to-back across two tiers in the reachable side rack.

\begin{center}
\begin{tikzpicture}[x=1cm,y=1cm,font=\small]
\foreach \x/\y/\lab/\rows in {0/1.25/{H / 16}/3,3.5/1.25/{H / 16}/3,0/0/{C / 8}/1,3.5/0/{O / 8}/1} {
\draw[fill=pale] (\x,\y) rectangle ++(3,1.08);
\node[anchor=west] at (\x+.12,\y+.89){\lab};
\foreach \i in {0,...,3} \foreach \j in {0,...,\rows} \draw[fill=accent!35] (\x+.7+\i*.5,\y+.15+\j*.17) circle (.055);
}
\end{tikzpicture}
\par\small Figure 2. Four baskets, two tiers; every dot represents one atom.
\end{center}


The 32 H permit eight methane molecules or four ethanes with hydrogen left over. Oxygen supports water and the hydrogen/oxygen reaction. These capacities are alternatives, not a promise that all maximum quantities can coexist. Recycling enables continued practice.

Selecting a recipe moves a seed from the shelf into the assembly area. Switching an unfinished recipe returns its atoms. A completed molecule can be selected and recycled; recycling either radical returns both fragments, preventing an orphaned rejoin state. Loose atoms can be returned explicitly and are recovered if dropped outside the working area. Reset restores all 48 atoms and cancels outstanding animation.

\begin{center}
\begin{tikzpicture}[node distance=7mm]
\node[state,text width=24mm] (s) {Shelf};\node[state,text width=24mm,right=of s] (l) {Loose / held};\node[state,text width=24mm,right=of l] (m) {Molecule};\node[state,text width=24mm,right=of m] (r) {Fragments / products};
\draw[flow] (s)--(l);\draw[flow] (l)--(m);\draw[flow] (m)--(r);
\draw[flow] (r.south)--++(0,-.55)-| node[pos=.3,below,font=\small]{Return / recycle / reset} (s.south);
\end{tikzpicture}
\par\small Figure 3. Ownership changes; total element counts remain constant.
\end{center}


Every token has a persistent identity and one owner. Reaction and fragment operations transfer identities rather than creating extra stock. During reaction animation, tokens remain owned by inputs until reassigned to products. No atom is silently discarded to satisfy an object-count limit.
